\documentclass[12pt,a4paper]{article}

% ── Packages ──────────────────────────────────────────────────────────────────
\usepackage[margin=2.5cm]{geometry}
\usepackage{graphicx}
\usepackage{hyperref}
\usepackage{xcolor}
\usepackage{listings}
\usepackage{enumitem}
\usepackage{titlesec}
\usepackage{fancyhdr}
\usepackage{tcolorbox}
\usepackage[T1]{fontenc}
\usepackage{lmodern}
\usepackage{longtable}
\usepackage{booktabs}
\usepackage{array}
\usepackage{soul}

% ── Colors ────────────────────────────────────────────────────────────────────
\definecolor{codebg}{HTML}{F5F5F5}
\definecolor{darkbg}{HTML}{1E293B}
\definecolor{accent}{HTML}{B91C1C}
\definecolor{green}{HTML}{166534}
\definecolor{codeblue}{HTML}{1E40AF}

% ── Code listing style ────────────────────────────────────────────────────────
\lstset{
  basicstyle=\ttfamily\small,
  backgroundcolor=\color{codebg},
  frame=single,
  rulecolor=\color{gray!40},
  breaklines=true,
  breakatwhitespace=true,
  tabsize=2,
  showstringspaces=false,
  columns=fullflexible,
  xleftmargin=1em,
  framexleftmargin=0.5em,
  keywordstyle=\color{codeblue}\bfseries,
  commentstyle=\color{gray},
  language=Go,
}

% ── Header / Footer ──────────────────────────────────────────────────────────
\pagestyle{fancy}
\fancyhf{}
\fancyhead[L]{\textsl{K8s Adaptive Workflows}}
\fancyhead[R]{\textsl{Issue Rebuttal Document}}
\fancyfoot[C]{\thepage}

% ── Tip/Warning boxes ────────────────────────────────────────────────────────
\newtcolorbox{srsbox}{colback=blue!5,colframe=codeblue,title=\textbf{What My SRS Says},fonttitle=\bfseries}
\newtcolorbox{codebox}{colback=green!3,colframe=green,title=\textbf{What My Code Actually Does},fonttitle=\bfseries}
\newtcolorbox{verdictbox}{colback=red!5,colframe=accent,title=\textbf{Verdict},fonttitle=\bfseries}
\newtcolorbox{factbox}{colback=yellow!5,colframe=yellow!60!black,title=\textbf{Key Fact},fonttitle=\bfseries}

% ══════════════════════════════════════════════════════════════════════════════
\begin{document}

% ── Title Page ────────────────────────────────────────────────────────────────
\begin{titlepage}
\centering
\vspace*{2cm}
{\Huge\bfseries Issue Rebuttal Document}\\[0.8cm]
{\LARGE K8s Adaptive Workflows}\\[0.5cm]
{\large Response to GitHub Issues \#2, \#3, and \#4}\\[2cm]

{\large
\begin{tabular}{ll}
\textbf{Repository:} & \url{https://github.com/Aman-Git-hala/k8s-adaptive-workflows} \\[0.3cm]
\textbf{SRS Reference:} & \texttt{k8s-adaptive-SRS-final.pdf} (25 pages, IEEE 830-1998) \\[0.3cm]
\textbf{Issues Filed By:} & \texttt{anshika289} (all three, April 16, 2026) \\[0.3cm]
\end{tabular}
}

\vspace{2cm}

{\large Prepared by:}\\[0.3cm]
{\large
Aman Githala (24293916037) \\
Ravi Kumar (24293916042) \\
Chandan Chaurasia (24293916019) \\
Ramkrishan Jangir (24293916033)
}

\vspace{2cm}

{\large Presented to:}\\[0.2cm]
{\large\textbf{Asst. Prof. Rekha R}}\\[0.2cm]
{\normalsize Department of Computer Science \& Engineering}

\vfill
{\large April 2026}
\end{titlepage}

\tableofcontents
\newpage

% ══════════════════════════════════════════════════════════════════════════════
\section{Introduction}
% ══════════════════════════════════════════════════════════════════════════════

On April 16, 2026, three issues were opened on my GitHub repository by the user \texttt{anshika289}. All three were filed within a short time window and follow the exact same template format.

After going through every single claim in these issues and comparing them line-by-line against my actual source code and my SRS document, I am confident that:

\begin{itemize}
  \item \textbf{Issue \#3} (marked ``Critical'') is \textbf{factually wrong}. The feature they say is missing is literally written in my code with clear comments. They did not read my source code.
  \item \textbf{Issue \#2} has a \textbf{small valid point} buried under exaggerated claims. The system handles cycles safely, but I could show a better error message. That said, calling it a ``bug'' is a stretch --- my system never crashes or freezes from cycles.
  \item \textbf{Issue \#4} describes a \textbf{feature that is not even in my SRS}. The data they want is already available through standard Kubernetes commands. Their framing makes it sound like a defect, but it is at best a nice-to-have improvement.
\end{itemize}

\subsection{A Note On Issue Quality}

Before I begin the detailed rebuttal, I want to point out some serious quality concerns with these issues:

\begin{enumerate}
  \item \textbf{Not a single line of code is referenced.} Across all three issues, there is no file path, no function name, no line number. A proper bug report on a codebase should reference the actual code.
  \item \textbf{No reproduction evidence.} None of the issues include logs, \texttt{kubectl} output, screenshots, or any proof that the described behavior was actually observed.
  \item \textbf{Identical template structure.} All three issues use the exact same format (Bug Type $\to$ Description $\to$ Steps $\to$ Expected $\to$ Actual $\to$ Impact $\to$ Fix), which strongly suggests they were written in a batch without individually investigating each claim.
  \item \textbf{The severity classifications are wrong.} Issue \#3 is marked ``Critical'' for a feature that is already fully implemented. This is either careless or intentionally misleading.
\end{enumerate}

\subsection{SRS Terminology (For Reference)}

My SRS document follows \textbf{IEEE 830-1998} conventions. The exact wording from SRS Section 1.2:

\begin{quote}
\textbf{MUST / SHALL:} Indicates a mandatory requirement. Failure to meet this requirement implies the system is non-compliant.

\textbf{SHOULD:} Indicates a recommended requirement that allows for flexibility during implementation.

\textbf{MAY:} Indicates an optional feature that will be implemented if time permits.
\end{quote}

This distinction matters a lot. A ``MUST'' requirement that is unmet is a real problem. But a feature that the SRS never mentions is not a ``bug'' --- it is a suggestion.

\newpage

% ══════════════════════════════════════════════════════════════════════════════
\section{Issue \#3 --- ``Resource Limit Enforcement Unclear (Critical)''}
% ══════════════════════════════════════════════════════════════════════════════

\textit{I am starting with Issue \#3 because it is marked ``Critical'' and is the most egregiously wrong.}

\subsection{What They Claim}

\begin{quote}
``Resource limits (maxResources) are defined but enforcement is not clearly guaranteed. [\ldots] verbally stated maxResource is 52 but no visible confirmation or logs proving enforcement. It takes unknown amount of time if number of resource is exceeded.''

\textbf{Impact:} ``May lead to pod eviction, OOM, or cluster instability.''
\end{quote}

\subsection{What My SRS Says}

\begin{srsbox}
\textbf{REQ-5.1} (SRS Section 4.5.3, Page 15): ``The system \textbf{MUST NOT} schedule a task if it would cause the cluster to exceed its resource quotas.''

\textbf{SRS Section 2.3} (Page 8): ``Resource Governance: The system strictly manages CPU and RAM usage to ensure it can run on low-cost hardware (`Free Tier') without crashing.''
\end{srsbox}

\subsection{What My Code Actually Does}

\begin{codebox}
The resource enforcement is implemented as a \textbf{hard gate} in \texttt{internal/optimizer/greedy.go}. Here is the exact code:
\end{codebox}

\textbf{Step 1:} Calculate how much CPU and memory is already being used by running tasks (lines 48--63):

\begin{lstlisting}
// Calculate current resource usage from running tasks.
currentCPU := resource.Quantity{}
currentMem := resource.Quantity{}
for _, ts := range taskStatuses {
    if ts.Phase == TaskPhaseRunning {
        if req, ok := ts.AllocatedResources.Requests[
            corev1.ResourceCPU]; ok {
            currentCPU.Add(req)
        }
        if req, ok := ts.AllocatedResources.Requests[
            corev1.ResourceMemory]; ok {
            currentMem.Add(req)
        }
    }
}
maxCPU := spec.MaxResources[corev1.ResourceCPU]
maxMem := spec.MaxResources[corev1.ResourceMemory]
\end{lstlisting}

\textbf{Step 2:} For every task that is ready to run, check if adding it would go over the limit. If yes, \textbf{skip it}. Do not schedule it. Wait for other tasks to finish first (lines 90--105):

\begin{lstlisting}
// Check if adding this task would exceed MaxResources.
taskCPU := taskResources.Requests[corev1.ResourceCPU]
taskMem := taskResources.Requests[corev1.ResourceMemory]

newCPU := currentCPU.DeepCopy()
newCPU.Add(taskCPU)
newMem := currentMem.DeepCopy()
newMem.Add(taskMem)

// If MaxResources is set, enforce the constraint.
if !maxCPU.IsZero() && newCPU.Cmp(maxCPU) > 0 {
    continue // Would exceed CPU limit, skip for now.
}
if !maxMem.IsZero() && newMem.Cmp(maxMem) > 0 {
    continue // Would exceed Memory limit, skip for now.
}
\end{lstlisting}

This is not hidden or complicated code. It is right there in the file, with clear comments that say \texttt{"Would exceed CPU limit, skip for now."}

\textbf{Step 3:} The controller also writes the \textbf{current resource usage} to the CRD status on every reconciliation cycle (controller line 249):

\begin{lstlisting}
wf.Status.CurrentResourceUsage =
    r.computeResourceUsage(wf.Status.TaskStatuses)
\end{lstlisting}

This is visible to anyone by running:
\begin{lstlisting}[language=bash]
kubectl get adaptiveworkflow <name> -o yaml
\end{lstlisting}

\subsection{About Their ``Unknown Amount of Time'' Complaint}

They write: ``It takes unknown amount of time if number of resource is exceeded.''

This is not a bug. This is \textbf{the entire point of my system}. When all the available resource budget is used up by running tasks, the optimizer waits for some tasks to finish before starting new ones. The ``unknown time'' is just the time it takes for pods to complete their work. That is how every single DAG scheduler in existence works --- Argo Workflows, Apache Airflow, Tekton, all of them.

My SRS calls this ``Bin Packing'' (Section 1.4): fitting as many tasks as possible onto available hardware without crashing the nodes. That is exactly what the code does.

\subsection{About Their ``Pod Eviction, OOM'' Concern}

They say the system ``may lead to pod eviction, OOM, or cluster instability.'' This is the \textbf{opposite} of what happens. The hard gate at lines 90--105 \textbf{prevents} the exact problems they are worried about. The system refuses to schedule a task if it would push resource usage over the limit. That is REQ-5.1 implemented word for word.

\begin{verdictbox}
\textbf{Issue \#3 is factually incorrect.} REQ-5.1 is fully implemented. The reviewer clearly did not open the file \texttt{internal/optimizer/greedy.go}. Marking this as ``Critical'' without reading the code is irresponsible. Their suggested fix (``Add resource accounting logs and enforcement validation'') already exists as the \texttt{CurrentResourceUsage} field in the CRD status.
\end{verdictbox}

\newpage

% ══════════════════════════════════════════════════════════════════════════════
\section{Issue \#2 --- ``No DAG Cycle Detection''}
% ══════════════════════════════════════════════════════════════════════════════

\subsection{What They Claim}

\begin{quote}
``The system does not visibly validate cyclic dependencies in DAG workflows. [\ldots] May cause infinite loops or scheduler freeze.''

Steps to reproduce: ``Create a cyclic DAG: A $\to$ B $\to$ C $\to$ A. Submit using CLI.''
\end{quote}

\subsection{What My SRS Says}

\begin{srsbox}
\textbf{REQ-1.2} (SRS Section 4.1.3, Page 12): ``The system \textbf{MUST} validate that the graph contains no cycles (infinite loops).''

\textbf{SRS Section 2.3} (Page 8): ``Workflow Ingestion \& Validation: The system accepts workflow definitions, validates them for logical correctness (e.g., no infinite loops), and prepares them for execution.''
\end{srsbox}

I will be honest here. My SRS does say ``MUST validate'' for cycle detection. So on paper, this is a valid requirement to bring up. \textbf{However}, the issue dramatically exaggerates the impact and makes claims that are physically impossible given my architecture.

\subsection{Why ``Infinite Loops or Scheduler Freeze'' Cannot Happen}

Let me walk through what actually happens if someone submits a cycle like A $\to$ B $\to$ C $\to$ A.

My optimizer (\texttt{internal/optimizer/greedy.go}, lines 72--83) only schedules a task if \textbf{ALL of its dependencies have succeeded}:

\begin{lstlisting}
// Check all dependencies are Succeeded.
ready := true
for _, dep := range task.Dependencies {
    depStatus, exists := taskStatuses[dep]
    if !exists || depStatus.Phase != TaskPhaseSucceeded {
        ready = false
        break
    }
}
if !ready {
    continue
}
\end{lstlisting}

Now think about what happens with a cycle A $\to$ B $\to$ C $\to$ A:
\begin{itemize}
  \item Task A is waiting for Task C to succeed.
  \item Task C is waiting for Task B to succeed.
  \item Task B is waiting for Task A to succeed.
  \item Nobody can ever start. All three stay in ``Pending'' forever.
\end{itemize}

\begin{factbox}
A cycle does \textbf{not} cause an infinite loop. It does \textbf{not} freeze the scheduler. It does \textbf{not} crash anything. The controller keeps reconciling every 10 seconds, checks if any task is ready, finds none are ready, and goes back to sleep. The workflow just sits there doing nothing. No CPU is wasted. No memory is leaked. No pods are created.
\end{factbox}

This is exactly the same behavior as Argo Workflows, by the way. If you submit a cyclic DAG to Argo, it does not crash either --- the stuck tasks just never progress.

\subsection{What IS Missing (And I Acknowledge It)}

What is missing is an \textbf{explicit error message} at submission time. Ideally, when a user submits a YAML with a cycle, the system should immediately reject it and say ``Error: your DAG contains a cycle involving tasks A, B, C.''

I do not have that validation step yet. The system handles cycles \textit{safely} (nothing breaks), but it does not handle them \textit{helpfully} (the user gets no clear message about why nothing is happening).

So REQ-1.2 is partially unmet --- not because the system is unsafe, but because it does not give a nice error. That is a usability enhancement, not a critical bug.

\subsection{Problems With Their Issue}

\begin{enumerate}
  \item They claim ``may cause infinite loops or scheduler freeze.'' This is \textbf{impossible}. My reconcile loop is stateless and runs on a fixed 10-second timer (controller line 264). It cannot loop infinitely or freeze.
  \item They provide no evidence that they actually tried submitting a cyclic DAG and observed a crash or freeze.
  \item Their ``Steps to Reproduce'' say ``Create a cyclic DAG: A $\to$ B $\to$ C $\to$ A. Submit using CLI.'' But they do not include any actual YAML, no CLI output, no logs. Did they actually try this?
\end{enumerate}

\begin{verdictbox}
\textbf{Partially valid, but significantly overstated.} The SRS says MUST validate cycles, and I do not have an explicit validation message. But the claim of ``infinite loops or scheduler freeze'' is completely wrong --- the architecture makes that impossible. This should be classified as a \textbf{Low severity enhancement request}, not a bug.
\end{verdictbox}

\newpage

% ══════════════════════════════════════════════════════════════════════════════
\section{Issue \#4 --- ``ML Predictions Not Visible (Usability)''}
% ══════════════════════════════════════════════════════════════════════════════

\subsection{What They Claim}

\begin{quote}
``The system uses ML for resource prediction but does not expose predictions to the user. You CANNOT see: Predicted CPU (e.g., 500m), Predicted memory (e.g., 256Mi), Predicted duration and Confidence/accuracy.''

\textbf{Impact:} ``User cannot verify or trust ML component.''
\end{quote}

\subsection{What My SRS Says}

\begin{srsbox}
Here are the ML-related requirements from my SRS (Section 4.3.3, Page 14):

\textbf{REQ-3.1:} ``The system SHALL record the actual resource usage of every completed task.''

\textbf{REQ-3.2:} ``The system SHALL employ a supervised regression model to estimate resource usage based on task metadata.''

\textbf{REQ-3.3:} ``Cold Start Strategy: If fewer than 10 historical execution records exist for a Task ID, the system MUST fallback to the static resource requests defined in the YAML.''

\textbf{REQ-3.4:} ``Training Constraint: Model retraining SHALL occur offline.''

\textbf{REQ-3.5:} ``If the model predicts a high chance of failure (OOM), the scheduler SHOULD proactively request more memory.''

\textbf{Notice what is NOT in the SRS}: There is no REQ-3.x that says ``the system MUST display predictions to the user.'' Not a single one. The SRS requirements focus on the ML pipeline itself --- record data, train model, predict resources, handle cold start. User-facing display of predictions was never a specified requirement.
\end{srsbox}

\subsection{But Actually, Predictions ARE Visible}

Even though the SRS does not require it, predictions are already visible through multiple channels:

\begin{enumerate}
  \item \textbf{CRD Status (kubectl):} Every task's allocated resources are written to the AdaptiveWorkflow status. You can see them by running:
  \begin{lstlisting}[language=bash]
kubectl get adaptiveworkflow <name> -o yaml
  \end{lstlisting}
  Each task in the status section shows its \texttt{allocatedResources.requests.cpu} and \texttt{allocatedResources.requests.memory}. These ARE the predicted/resolved values that the optimizer used.

  \item \textbf{Inference Engine Logs:} When running with Docker Compose, the Python inference service logs every single prediction to stdout:
  \begin{lstlisting}[language=bash]
docker compose logs inference-engine
  \end{lstlisting}

  \item \textbf{Controller Logs:} Every pod creation is logged with the task name (controller line 243):
  \begin{lstlisting}
log.Info("Created pod for task",
    "task", alloc.TaskName, "pod", ts.PodName)
  \end{lstlisting}

  \item \textbf{Web Dashboard:} The simulation view shows an adaptive vs.\ baseline comparison, which directly visualizes how the ML predictions affect scheduling.

  \item \textbf{PostgreSQL StateDB:} All completed task metrics (actual CPU, memory, duration) are recorded and can be queried:
  \begin{lstlisting}[language=SQL]
SELECT * FROM task_executions
    ORDER BY created_at DESC;
  \end{lstlisting}
\end{enumerate}

\subsection{About Their ``Confidence/Accuracy'' Demand}

They say users should be able to see ``Confidence/accuracy'' for each prediction. Let me be real here:

\begin{itemize}
  \item \textbf{No production workflow system does this.} Not Argo Workflows. Not Apache Airflow. Not Tekton. Not Prefect. Not Dagster. None of them show ``confidence scores'' for resource estimates in their CLI output.
  \item Per-prediction confidence is a concept from ML model evaluation, not from runtime scheduling systems. It makes sense in a Jupyter notebook, not in a \texttt{kubectl} command.
  \item Even if I wanted to add it, my SRS does not require it.
\end{itemize}

\begin{verdictbox}
\textbf{Issue \#4 is a feature request, not a bug.} The SRS contains zero requirements about user-facing prediction display. Predictions are already visible through \texttt{kubectl}, inference logs, and the dashboard. Their claim that ``You CANNOT see'' predictions is incorrect. The ``Confidence/accuracy'' demand is unrealistic and not standard in any competing system. Severity: \textbf{Low / Enhancement at best}.
\end{verdictbox}

\newpage

% ══════════════════════════════════════════════════════════════════════════════
\section{Summary}
% ══════════════════════════════════════════════════════════════════════════════

\begin{longtable}{|p{2.2cm}|p{2cm}|p{3.5cm}|p{2.2cm}|p{3cm}|}
\hline
\textbf{Issue} & \textbf{Claimed Severity} & \textbf{SRS Requirement} & \textbf{Met?} & \textbf{My Assessment} \\
\hline
\#2 DAG Cycle Detection & Usability (Bug) & REQ-1.2: MUST validate no cycles & Partially (no error msg, but cycle-safe) & Low / Enhancement \\
\hline
\#3 Resource Enforcement & \textbf{Critical} (Bug) & REQ-5.1: MUST NOT exceed quotas & \textbf{Fully met} (greedy.go L90--105) & \textbf{Invalid} --- code not reviewed \\
\hline
\#4 ML Prediction Visibility & Usability (Bug) & No SRS requirement exists for this & Visible via \texttt{kubectl}, logs, dashboard & Low / Enhancement \\
\hline
\end{longtable}

\subsection{Overall Assessment}

Out of 3 issues:
\begin{itemize}
  \item \textbf{1 is flat-out wrong} (Issue \#3 --- the feature exists and is implemented)
  \item \textbf{1 has a small valid point} but with wildly exaggerated impact claims (Issue \#2)
  \item \textbf{1 is a feature request} for something not in the SRS, framed as a bug (Issue \#4)
\end{itemize}

\subsection{What Proper Issues Would Look Like}

For comparison, a well-written issue on a codebase like mine would include:
\begin{enumerate}
  \item A reference to the specific file and line number where the problem is.
  \item Actual reproduction steps with real YAML files, not vague instructions.
  \item Actual output --- what command was run, what the terminal showed, what went wrong.
  \item A clear distinction between ``the code crashes'' (bug) and ``I wish the code did X'' (feature request).
\end{enumerate}

None of the three issues meet any of these standards. They read like they were written based on a quick scan of my README, not from actually running the code or reading the implementation.

\newpage

% ══════════════════════════════════════════════════════════════════════════════
\section{Code References Quick Index}
% ══════════════════════════════════════════════════════════════════════════════

For easy verification, here is exactly where to look for each claimed ``missing'' feature:

\begin{longtable}{|p{4cm}|p{5cm}|p{4cm}|}
\hline
\textbf{What They Claim is Missing} & \textbf{File} & \textbf{Lines} \\
\hline
Resource limit enforcement & \texttt{internal/optimizer/greedy.go} & Lines 48--63 (usage sum), Lines 90--105 (hard gate) \\
\hline
Resource usage tracking & \texttt{internal/controller/adaptiveworkflow\_controller.go} & Line 249: \texttt{computeResourceUsage()} \\
\hline
Resource usage in CRD & \texttt{api/v1/adaptiveworkflow\_types.go} & Line 100: \texttt{CurrentResourceUsage} field \\
\hline
Dependency checking (cycle safety) & \texttt{internal/optimizer/greedy.go} & Lines 72--83 (all deps must be Succeeded) \\
\hline
ML predictions in status & \texttt{internal/controller/adaptiveworkflow\_controller.go} & Lines 235--237: \texttt{AllocatedResources} written per task \\
\hline
Inference engine predictions & \texttt{inference-engine/server.py} & \texttt{Predict()} RPC handler --- logs all predictions \\
\hline
Metrics to StateDB & \texttt{internal/controller/adaptiveworkflow\_controller.go} & Lines 312--340: \texttt{r.StateDB.RecordExecution()} \\
\hline
\end{longtable}

\vspace{1cm}

\begin{center}
\rule{0.8\textwidth}{0.5pt}

\vspace{0.5cm}
\textit{This document was prepared in response to issues filed on my repository. I am happy to demonstrate any of the above code paths live during the evaluation.}

\vspace{0.3cm}
--- Aman Githala, Ravi Kumar, Chandan Chaurasia, Ramkrishan Jangir
\end{center}

\end{document}
